\documentclass[12pt,a4paper]{article}
\usepackage[T2A]{fontenc}
\usepackage[utf8]{inputenc}
\usepackage[russian]{babel}
\usepackage{geometry}
\usepackage{enumitem}
\usepackage{listings}
\usepackage{xcolor}
\usepackage{hyperref}
\usepackage{longtable}
\usepackage{array}
\geometry{left=25mm,right=15mm,top=20mm,bottom=20mm}
\setlength{\parindent}{1.25cm}
\setlength{\parskip}{0pt}
\setlength{\emergencystretch}{4em}
\lstset{basicstyle=\ttfamily\small,frame=single,breaklines=true,
showstringspaces=false,columns=fullflexible,keepspaces=true}
\hypersetup{colorlinks=true,linkcolor=black,urlcolor=blue}
\begin{document}
\begin{titlepage}
\begin{center}
Федеральное государственное автономное образовательное учреждение\\
высшего образования\\[2mm]
\textbf{«Национальный исследовательский университет\\
``Высшая школа экономики''»}\\[3mm]
Факультет компьютерных наук\\
Департамент программной инженерии
\vfill
{\Large\bfseries Отчёт по индивидуальному заданию № 4}\\[4mm]
{\large по дисциплине «Архитектура компьютера и операционные системы»}\\[7mm]
{\Large\bfseries Вариант 14. «Задача об умных цветах»}\\[3mm]
{\large Имитационное моделирование с использованием TCP}
\vfill
\begin{flushright}
Выполнила: Гончарова Екатерина\\
Группа: БПИ-253
\end{flushright}
\vfill
Москва, 2026
\end{center}
\end{titlepage}
\tableofcontents
\newpage

\section{Постановка задачи}
В индивидуальном задании требуется разработать несколько консольных приложений для Linux, которые взаимодействуют по сети с использованием протокола TCP. Вариант 14 посвящён моделированию умных цветов и садовников. Цветы постепенно теряют влагу, а при достижении заданного порога формируют заявки на полив. Садовники выбирают доступные заявки, подходят к цветкам и выполняют полив. Если цветок пересох до критического значения, он погибает. Повторный полив во время защитного интервала также может привести к гибели.

В данной реализации участники обмениваются данными посредством TCP-соединений. Сервер и клиенты могут работать на одном компьютере или на нескольких узлах локальной сети. Предусмотрены наблюдение за состоянием сада, повторные подключения клиентов и согласованное завершение приложений.

\section{Модель и входные данные}
\subsection{Цветы и садовники}
Для представления объектов используются структуры \texttt{Flower}, \texttt{Gardener} и \texttt{Settings}. Для цветка определены состояния \texttt{NORMAL}, \texttt{WILTING}, \texttt{WATERING}, \texttt{RECOVERING}, \texttt{DRY} и \texttt{DEAD\_}\allowbreak\texttt{OVERWATERED}. Начальная влажность равна 100. На очередном шаге она уменьшается на индивидуальную скорость высыхания, пока цветок не находится в состоянии полива или восстановления. Заявка возникает при достижении порога запроса, а гибель от засыхания --- при достижении нижнего порога.

Садовник принимает только свободную заявку. Используются две стратегии: \texttt{first} --- первая доступная заявка, \texttt{driest} --- заявка от цветка с наименьшей влажностью. После назначения садовник сначала выполняет работу по перемещению, затем поливает цветок в течение установленного числа тактов. Успешный полив восстанавливает влажность до 100. После полива действуют защитный интервал и период восстановления.

\subsection{Настройки симуляции}
\begin{longtable}{|p{0.39\textwidth}|p{0.49\textwidth}|}
\hline
Параметр & Допустимые значения\\ \hline
\endhead
\texttt{flower\_count} & От 1 до 100\\ \hline
\texttt{drying\_rates} & По одному положительному числу на цветок\\ \hline
\texttt{request\_threshold} & От 1 до 100\\ \hline
\texttt{dry\_threshold} & От 0 до 99, строго меньше порога заявки\\ \hline
\texttt{watering\_duration} & Положительное число тактов\\ \hline
\texttt{recovery\_duration} & Неотрицательное число тактов\\ \hline
\texttt{protection\_interval} & Неотрицательное число тактов\\ \hline
\texttt{gardener\_count} & От 1 до 10\\ \hline
\texttt{gardener\_speed} & Положительное целое число\\ \hline
\texttt{request\_policy} & \texttt{first} или \texttt{driest}\\ \hline
\end{longtable}

Моделирование длится не более 50 тактов. Обычная задержка между тактами составляет 120~мс; её можно изменить переменной окружения \texttt{SMART\_FLOWERS\_TICK\_MS}. Этот параметр влияет только на скорость вывода событий. Входные данные можно получить из конфигурационного файла либо ввести в терминале сервера.

\section{Архитектура приложения}
\subsection{Состав программ}
Реализация состоит из четырёх исполняемых программ:
\begin{itemize}[nosep]
\item \texttt{server} --- хранит актуальное состояние сада, принимает TCP-подключения, запускает модельные такты и формирует журнал событий;
\item \texttt{flower} --- клиент конкретного цветка, получает команды такта и его локальное состояние;
\item \texttt{gardener} --- клиент конкретного садовника, получает команды такта и информацию о назначении;
\item \texttt{observer} --- наблюдатель, получающий сводную информацию обо всём саде.
\end{itemize}

Сервер хранит текущее состояние всех объектов, но не вычисляет за них изменения. Клиенты не имеют доступа к памяти сервера: цветок получает собственное состояние и пороги, вызывает \texttt{update\_flower()} и отправляет результат. Садовник получает своё состояние, настройки полива и снимок цветов, вызывает \texttt{update\_gardener()} и возвращает изменения. Вычисления выполняются независимыми сетевыми приложениями, а сервер согласует результаты.

Для каждого соединения создаётся обслуживающий поток \texttt{pthread}; общие данные защищает \texttt{pthread\_mutex\_t}. Модельный поток последовательно запрашивает результаты сначала у цветов, затем у садовников. Сервер проверяет номер такта и идентификатор клиента, применяет полученные изменения и публикует новое состояние для наблюдателей.

До подключения всех необходимых цветов и садовников симуляция не начинается. Если рабочий клиент отключился, сервер приостанавливается на текущем шаге и ждёт повторного подключения того же участника. Состояние не теряется, а пропущенные действия сервер самостоятельно не выполняет. Таким образом, повторное подключение не требует повторной инициализации состояния сада.

\subsection{Структура исходного кода}
Модель сада реализована в модулях \texttt{model.c}, \texttt{flower\_logic.c}, \texttt{gardener\_logic.c}, \texttt{request\_logic.c} и \texttt{statistics.c}. Ввод и проверку параметров обеспечивают \texttt{input.c}, \texttt{config\_parser.c} и \texttt{input\_validation.c}. Сетевые соединения обрабатываются в \texttt{network.c}, \texttt{server.c} и \texttt{client.c}. Модуль \texttt{worker\_protocol.c} преобразует состояния объектов в текстовые сообщения и обратно. Регистрацией клиентов и приёмом их ответов занимается \texttt{server\_clients.c}. Журналирование защищено от одновременной записи несколькими потоками.

\section{Передача данных по TCP}
\subsection{Создание соединений}
Сервер вызывает \texttt{socket()}, \texttt{bind()}, \texttt{listen()} и \texttt{accept()}. IP-адрес привязки и порт указываются при запуске в командной строке. Клиент вызывает \texttt{socket()} и \texttt{connect()}, указывая адрес сервера. Это позволяет использовать программу как на одном компьютере, так и в локальной сети. Для работы на удалённом компьютере необходимо обеспечить сетевую доступность выбранного TCP-порта.

Команды представлены строками текста, заканчивающимися символом перевода строки. Такой формат прост для отладки и позволяет просматривать сообщения обычными средствами терминала. Чтение ограничено размером буфера; слишком длинные сообщения не обрабатываются как корректные команды.

\subsection{Основные сообщения}
\begin{longtable}{|p{0.34\textwidth}|p{0.54\textwidth}|}
\hline
Сообщение & Назначение\\ \hline
\endhead
\texttt{HELLO FLOWER 1} & Регистрация цветка с номером 1\\ \hline
\texttt{HELLO GARDENER 1} & Регистрация садовника с номером 1\\ \hline
\texttt{HELLO OBSERVER} & Подключение наблюдателя\\ \hline
\texttt{WELCOME} & Подтверждение регистрации\\ \hline
\texttt{FJOB ...} & Состояние цветка и параметры для вычисления его шага\\ \hline
\texttt{FRESULT ...} & Рассчитанное клиентом состояние цветка\\ \hline
\texttt{GJOB ...} & Состояние садовника и снимок сада\\ \hline
\texttt{GRESULT ...} & Результат работы садовника и изменение цветка\\ \hline
\texttt{STATE ...} & Текущее состояние объекта для его клиента\\ \hline
\texttt{SNAPSHOT ...} & Сводка для наблюдателя\\ \hline
\texttt{FLOWER ...}, \texttt{GARDENER ...} & Данные объектов внутри снимка\\ \hline
\texttt{END\_SNAPSHOT} & Завершение снимка\\ \hline
\texttt{FINISH ...} & Уведомление о завершении симуляции\\ \hline
\end{longtable}

Например, клиент цветка отправляет \texttt{HELLO FLOWER 1} и получает \texttt{WELCOME}. После регистрации остальных работников сервер передаёт сообщение \texttt{FJOB}, содержащее номер такта, пороги и поля цветка. Клиент вычисляет новое состояние с помощью \texttt{update\_flower()} и отвечает \texttt{FRESULT}. Аналогично садовник рассчитывает свой шаг по \texttt{GJOB} и отправляет \texttt{GRESULT}.

Сервер применяет результаты только для ожидающего работника и нужного такта, поэтому повторный ответ не должен приводить к повторному поливу. Если соединение обрывается до получения результата, шаг выполняется заново после переподключения, исходя из последнего сохранённого сервером состояния. После обработки всех цветов и садовников сервер рассылает обновлённые снимки наблюдателям.

\section{Варианты отображения информации}
\subsection{Раздельное отображение (4--5 баллов)}
Клиент цветка выводит события своего шага, клиент садовника --- действия по принятию заявки и поливу. Для выбора заявки садовник получает снимок состояний цветов, но не изменяет их напрямую: результат передаётся серверу сообщением. Сервер ведёт общий технический журнал; сводное отображение всех объектов предусмотрено отдельным наблюдателем.

\subsection{Общий наблюдатель (6--7 баллов)}
Программа \texttt{observer} получает полный снимок сада: номер такта, количество живых и погибших цветов, число поливов, текущую влажность и состояние каждого цветка, занятость и задания садовников. Поэтому за ходом всей симуляции можно следить без открытия терминалов отдельных работников.

\subsection{Несколько наблюдателей (8 баллов)}
Сервер обслуживает несколько TCP-подключений одновременно. Наблюдатели независимы: отключение одного не влияет на остальных, повторное подключение получает актуальный снимок. Встроено ограничение общего числа обслуживаемых соединений. При невозможности отправить данные медленному клиенту сервер прекращает его соединение, не ожидая бесконечно освобождения сетевого буфера.

\subsection{Повторные подключения и завершение клиентов}
Клиент цветка или садовника может завершиться и затем подключиться снова под тем же номером. Во время отсутствия нужного работника модельный поток ожидает его возвращения, не изменяя данный объект самостоятельно. Одновременная регистрация двух активных клиентов с одной ролью и номером отклоняется. При штатном завершении или обработке \texttt{SIGINT}/\texttt{SIGTERM} сервер отправляет подключённым клиентам \texttt{FINISH}, после чего они закрывают соединения. При обрыве сети без уведомления клиент также прекращает работу, но может быть запущен повторно.

\section{Пользовательский интерфейс и запуск}
\subsection{Сборка и демонстрация}
Все команды выполняются из корня проекта. Для сборки используется:
\begin{lstlisting}
make
\end{lstlisting}

Для автоматической демонстрации с выбранным набором параметров:
\begin{lstlisting}
bash scripts/run_demo.sh tests/balanced.cfg
bash scripts/run_demo.sh tests/overwatering.cfg 19015
\end{lstlisting}

Скрипт определяет количество цветов и садовников из файла, запускает сервер, работников и одного наблюдателя. Его задача --- упростить локальную демонстрацию, при этом все приложения остаются отдельными процессами и взаимодействуют только через TCP.

\subsection{Запуск из нескольких терминалов}
Для конфигурации с одним цветком используются следующие команды:
\begin{lstlisting}
# Terminal 1
./build/server 127.0.0.1 19014 tests/one_flower.cfg
# Terminal 2
./build/flower 127.0.0.1 19014 1
# Terminal 3
./build/gardener 127.0.0.1 19014 1
# Terminal 4
./build/observer 127.0.0.1 19014
\end{lstlisting}

Вместо файла параметров можно запустить сервер с двумя аргументами:
\begin{lstlisting}
./build/server 127.0.0.1 19014
\end{lstlisting}

Тогда программа последовательно запрашивает значения в терминале. Если введённое число нарушает ограничения, предлагается повторить ввод. Адрес сервера и порт задаются при запуске всех клиентов. Для распределённого запуска сервер может прослушивать \texttt{0.0.0.0}, а клиенты должны указывать доступный IP-адрес компьютера с сервером.

\section{Тестирование и результаты}
Для проверки реализован скрипт \texttt{tests/check\_tcp.py}, запускаемый командой \texttt{make check}. Он проверяет, что сервер ждёт подключения работников, затем выполняет модель, допускает переподключение цветка и наблюдателя и отклоняет повторную регистрацию номера. Для моделирования предусмотрено тринадцать конфигурационных сценариев с разными параметрами, включая сад из 100 цветов и 10 садовников.

\begin{center}
\begin{tabular}{|p{0.58\textwidth}|p{0.3\textwidth}|}
\hline
Проверка & Результат\\ \hline
Отсутствие подключённых работников & Модель не начинается\\ \hline
Переподключение цветка и двух наблюдателей & Пройдено\\ \hline
Повторная регистрация номера & Отклоняется\\ \hline
Тринадцать конфигурационных сценариев & Завершились успешно\\ \hline
100 цветов и 10 садовников по TCP & Завершились успешно\\ \hline
Остановка по \texttt{SIGINT} до старта & Сервер завершился\\ \hline
\end{tabular}
\end{center}

При одном из запусков \texttt{tests/one\_flower.cfg} после 50 тактов цветок оставался живым, влажность составляла 93, количество завершённых поливов --- 7. Точный результат определяется параметрами сценария и обработанными шагами работников. В журнале \texttt{simulation.log} сохраняются сообщения сервера о подключениях, тактах, принятых результатах и итоговом состоянии сада.

\section{Особенности и ограничения реализации}
Сервер хранит согласованное состояние сада и проверяет ответы работников. При этом действия цветов и садовников вычисляются клиентами, а не заменяются локальными шагами сервера. Такой подход позволяет продемонстрировать зависимость модели от сетевого взаимодействия и сохраняет возможность переподключения. Если обязательный клиент отсутствует, симуляция намеренно ждёт его подключения; завершить ожидание можно сигналом \texttt{SIGINT} или \texttt{SIGTERM}.

Для обмена используется строковый протокол TCP. Состояния передаются отдельными числовыми полями, чтобы не зависеть от внутреннего расположения структур C в памяти. Отправка и чтение данных имеют ограничение по времени; при отсутствии ответа сервер закрывает зависшее соединение и ожидает новое подключение работника. Из-за последовательного ожидания каждого шага скорость работы зависит от задержек сети, особенно при большом количестве цветов.

Приложение рассчитано на учебную локальную сеть. Шифрование и аутентификация клиентов не предусмотрены. Состояние не сохраняется между запусками сервера, а при его завершении все текущие TCP-сессии закрываются.

\section{Вывод}
Разработана клиент-серверная система для моделирования умных цветов. Она использует отдельные TCP-приложения для цветов, садовников и наблюдателей. Клиенты цветов и садовников самостоятельно рассчитывают изменения своей модели и отправляют результаты серверу по TCP. Реализованы общий наблюдатель, одновременное подключение нескольких наблюдателей, переподключение работников и уведомление об остановке сервера. Проверены обычные сценарии и запуск с большим количеством участников.

\end{document}
